\section{Design Specification}\label{sec:design}

The proposed forecasting framework is designed as an \textbf{eight-stage
processing pipeline}, as illustrated in \textbf{Figure 1}. The first
four stages are common to all four forecasting models, ensuring that
every model receives the same input data, feature set, and
training--testing split. This shared preprocessing guarantees a fair and
controlled comparison by eliminating differences caused by data
preparation.

After the common preprocessing stages, the pipeline branches into two
independent prediction paths. One branch generates predictions using the
LSTM-based machine learning model, while the other produces predictions
using the language model with or without retrieval support, depending on
the experimental configuration. The outputs from these two branches are
then combined through a calibrated fusion module to create the hybrid
prediction. Finally, the evaluation stage compares the performance of
all four forecasting models using the same set of predefined metrics.

The framework also includes a Flask-based web application that uses the
same implementation modules as the experimental pipeline. Since both the
research experiments and the user interface rely on the identical
codebase, the predictions and evaluation results displayed by the
application remain fully consistent with the results reported in this
study. This design improves reproducibility and ensures that the
experimental findings accurately represent the deployed system.

\textbf{Figure 1.} Overview of the proposed eight-stage forecasting
pipeline. The initial four stages perform common data processing for all
forecasting models, followed by separate prediction branches, a
calibrated fusion module, and a unified evaluation stage.

\begin{figure}[H]
\centering
\includegraphics[width=0.64\textwidth]{fig-pipeline.png}
\caption{The eight-stage pipeline. Stages 1-4 are shared by all four arms; stages 5 and 6 are the independent branches whose outputs stage 7 fuses.}
\label{fig:pipeline}
\end{figure}

\subsection{Prediction Models}\label{sec:prediction-models}

The proposed framework consists of four forecasting models that are
evaluated under identical experimental conditions. The machine learning
model is implemented using a compact LSTM architecture to reduce the
risk of overfitting, as each training fold contains only a limited
number of sequential samples. The network includes \textbf{two LSTM
layers with 64 and 32 hidden units}, followed by a \textbf{dropout layer
with a rate of 0.3}, a \textbf{fully connected layer containing 16
neurons}, and a \textbf{sigmoid output layer} that predicts the
probability of an upward price movement.

The LSTM architecture was selected after comparing it with four
alternative machine learning models: \textbf{Logistic Regression, Random
Forest, XGBoost, and GRU}. Each candidate model was trained using the
same walk-forward validation strategy, identical feature set, five-day
embargo period, and independently optimized decision threshold. The
comparative performance of these models is presented in \textbf{Section
6.1}, and the LSTM was chosen as the deterministic forecasting model for
the remainder of the study.

The language-model experiments are designed so that the only difference
between the two LLM-based systems is the availability of retrieved
external evidence. Both models use \textbf{Gemini 2.5 Flash} with the
same prompt structure, JSON response format, temperature setting, and
market information. The retrieval-enhanced model receives an additional
\textbf{EVIDENCE} section containing relevant historical information and
is instructed to use only this retrieved evidence when referring to
market events or external knowledge.

Maintaining identical configurations ensures that any observed
performance difference between the two language models can be attributed
solely to the retrieval component. In addition, each model is required
to report the numerical values and evidence identifiers used during
prediction. These outputs provide the information required for the
automated hallucination analysis described in \textbf{Section 6.4},
enabling the verification of numerical correctness and citation
accuracy.

\subsection{Fusion Algorithm}\label{sec:fusion-algorithm}

The hybrid forecasting model combines the outputs of the LSTM and the
retrieval-enhanced language model through a calibrated fusion strategy.
During the initial stage of development, a simple average of the
probability scores generated by the two models was tested. However, this
approach produced an overall accuracy of \textbf{53.7\%}, which was
lower than the performance of both the standalone LSTM
(\textbf{52.67\%}) and the grounded language model (\textbf{57.67\%}).
This result demonstrated that directly averaging probability values from
different models can reduce forecasting performance.

The primary reason for this behaviour is that the two models produce
probability estimates using different confidence scales. The LSTM
generates probabilities within a relatively narrow range and predicts an
upward movement only when its output exceeds approximately
\textbf{0.59}. In contrast, the language model uses a different
confidence distribution and applies an effective decision threshold
close to \textbf{0.60}. Since these probabilities are calibrated
differently, averaging them directly combines two incompatible
confidence scales and results in an unreliable decision boundary.

To overcome this limitation, each model output is first transformed into
a normalized confidence value representing its signed distance from its
own decision threshold. The hybrid prediction is then computed using the
following equations:

\begin{align*}
z_{\mathrm{arm}} &= \operatorname{logit}(p_{\mathrm{arm}})-\operatorname{logit}(t_{\mathrm{arm}}),\\
z_{\mathrm{hybrid}} &= (1-w)z_{\mathrm{ML}} + w z_{\mathrm{LLM\_RAG}},\\
p_{\mathrm{hybrid}} &= \operatorname{sigmoid}(z_{\mathrm{hybrid}}), \quad \text{decide UP if } p_{\mathrm{hybrid}} \geq 0.5.
\end{align*}

This transformation converts each model output into a common confidence
scale before fusion. A positive value indicates that the corresponding
model predicts an upward movement, while a value of zero represents its
own decision boundary. Since both models are centered around their
respective thresholds, the fused confidence score naturally uses
\textbf{0.5} as its classification boundary. Consequently, the hybrid
model does not require an additional decision threshold, reducing the
number of fitted parameters compared with the individual forecasting
models.

The fusion weight \textbf{(w)} is estimated separately for \textbf{CALM}
and \textbf{SHOCK} market conditions using the log-loss values obtained
from the independent tuning dataset. Instead of selecting weights that
maximize accuracy on the tuning data, the estimated values are partially
regularized toward an equal-weight combination. This strategy was
adopted for two reasons. First, the tuning dataset contains a different
distribution of market regimes compared with the final evaluation
dataset, and aggressively optimizing the weights could introduce bias
toward one market condition. Second, log-loss evaluates the complete
probability distribution rather than only the final classification,
making it a more reliable measure for estimating model confidence using
a relatively small dataset.

The resulting fusion weights are \textbf{0.492} for calm-market
conditions and \textbf{0.487} for shock-market conditions. Although
these values are close to an equal-weight combination, they are obtained
through a data-driven calibration process rather than being manually
assigned. This approach allows the hybrid model to adapt if one
forecasting model consistently demonstrates higher reliability than the
other while avoiding unnecessary overfitting to the tuning dataset.

\subsection{Personalization}\label{sec:personalization}

To provide recommendations that match different investment preferences,
the proposed system supports \textbf{three user risk profiles}:
\textbf{Cautious}, \textbf{Balanced}, and \textbf{Aggressive}. Each
profile applies a different minimum confidence threshold before the
system recommends taking action. The required confidence levels are
\textbf{70\%} for cautious users, \textbf{60\%} for balanced users, and
\textbf{55\%} for aggressive users.

It is important to note that the selected risk profile influences only
the investment recommendation presented to the user. The predicted
market direction and confidence score produced by the forecasting models
remain unchanged. This separation ensures that personalization affects
only the decision support provided to the user and does not influence
the forecasting models themselves.

Maintaining this distinction is essential for fair model evaluation.
Since the prediction process is independent of the recommendation
strategy, the reported forecasting accuracy reflects only the predictive
capability of the models. Similarly, the personalization score used in
the response quality evaluation measures the effectiveness of adapting
recommendations to different user preferences rather than artificially
improving the prediction results.
